\color{unimain}{\subsection{Защита от непереключения фаз}}
\color{black}

\begin{enumerate}
\item
Защита от непереключения фаз (ЗНФ) предназначена для обнаружения расхождения полюсов
выключателя с пофазным приводом при его включении или отключении. 

\item
Функциональная схема ЗНФ представлена на рисунке \ref{pict_znf:ZNF}.

\begin{figure}[H]
\centering
\includegraphics[width=0.8\textwidth,height=0.8\textheight,keepaspectratio]{\fbpath/08.01 СВ/661.3441 ЗНФ/img/ZNF.pdf}
\captionof{figure}{Функциональная схема ЗНФ}\label{pict_znf:ZNF}
\end{figure}

\item
Ввод в работу ЗНФ выключателя осуществляется с помощью программного ключа \textbf{«ЗНФ»}.
Факт расхождения полюсов может устанавливаться одним из следующих способов:
\begin{itemize}
	\item на основании внешнего сигнала от сборки блок-контактов выключателя;
	\item функцией фиксации положения выключателя (ФПВ), которая пофазно контролирует состояние блок-контактов выключателя и в случае расхождения полюсов формирует сигнал \textit{«Внутр. пуск. ЗНФ»}.
\end{itemize}

\item
В зависимости от используемого способа соответствующий сигнал должен назначаться на
конфигурируемый вход \textit{«ЗНФ СВ/ЗНФ:Пуск»}.

\item
Действие на отключение выключателя с запретом АПВ осуществляется с независимой выдержкой
времени \textbf{«Тср ЗНФ»}. Задаваемая выдержка времени предназначена для отстройки от разновременности
переключения блок-контактов выключателя.

\item
Предусмотрена возможность блокировки ЗНФ от внешнего сигнала \textit{«ЗНФ СВ/ЗНФ: Блокировка»} на время, которое
задается уставкой \textbf{«Тдеблок ЗНФ»}.

\item
Оперативный ввод/вывод ЗНФ производится с помощью БУ (блоков управления функциями) \textit{«ЗНФ СВ/ЗНФ: Вывод»}. 

\item
Уставочные параметры функции ЗНФ представлены в таблице \ref{app_set:znf}.

\end{enumerate}